\chapter{Conclusions and recommendations}
\label{ch:conclusions}

This chapter summarizes the final conclusions of the project and translates the
main findings into concrete recommendations for future work.

\section{Conclusions}

The objective of this project was to develop and validate an efficient and
scalable computer-vision workflow that enables a service-robot platform to detect
and classify human-interaction elements in hospital environments using embedded
hardware. Based on the results presented in this report, this objective was
achieved for the targeted elevator-interaction use case.

The project demonstrates that a practical embedded perception workflow can be
constructed from four main components: a strong instance-segmentation baseline
for button localization, a crop-based \gls{ocr} stage for button interpretation,
\gls{tensorrt} export for hardware-efficient inference, and a modular \gls{ros}
pipeline that integrates detection, 3D projection, text recognition, and
deployment data collection.

The main conclusions are as follows.

\begin{enumerate}
    \item Instance segmentation is an effective method for elevator-button perception because it supports both accurate 2D localization and mask-based depth estimation for 3D target generation.
    \item The baseline \gls{yolo}v8s segmentation model was the strongest deployment candidate among the evaluated detector variants. Hyperparameter tuning improved validation performance, but not the final held-out test ranking.
    \item Environment-specific fine-tuning with \gls{adc} data from Han Building R29 produced the most practically valuable \gls{yolo} improvement, substantially increasing target-domain performance while preserving original-domain performance. In the controlled single-button R29 experiment, the adapted detector produced more frames with detections (138 vs.\ 106) and longer mean and median valid-detection distances (0.4282~m vs.\ 0.3527~m and 0.4102~m vs.\ 0.3308~m, respectively) after the initial model's manually reviewed long-range false positives were excluded from the valid-distance comparison. The corresponding confidence values were very similar, indicating that the main practical gain was more reliable target coverage rather than a large confidence increase. A separate reviewed deployment-style university run further showed a clear practical trend: accepted \gls{yolo} detection confidence decreases as distance increases, which strengthens the case for treating distance as an important reliability cue during deployment.
    \item The selected stage-2 \gls{ocr} model improved held-out recognition quality and formed a suitable semantic interpretation stage for the full pipeline.
    \item \gls{tensorrt} export was beneficial for both deployed models. The effect was moderate but clear for \gls{yolo} inference time, and very strong for \gls{ocr} per-crop wall time.
    \item The final \gls{ros} pipeline achieved practical embedded runtime performance on the \gls{jetson} Orin Nano, with the best measured configuration operating at about 13.1~\gls{fps} (approximately 76~ms/frame) and the collector-enabled configuration remaining usable at about 10~\gls{fps} (approximately 100~ms/frame).
    \item Scalability is supported not only by model architecture choices, but also by the inclusion of the \gls{adc}, which enables structured deployment-time data collection for future retraining.
\end{enumerate}

Taken together, these conclusions answer the main research question: a practical
computer-vision workflow for elevator-button perception can be achieved on
embedded hardware by combining segmentation-based detection, task-specific
recognition, deployment-oriented optimization, and iterative deployment data
collection within a modular \gls{ros} architecture.

\section{Recommendations}

The following recommendations are proposed for future work and further
development of the system.

\begin{enumerate}
    \item Extend the perception scope beyond elevator buttons to additional hospital interaction elements such as access readers, door controls, and other panel interfaces.
    \item Perform larger and more controlled deployment trials in multiple buildings so that domain shift, lighting variation, and panel diversity can be evaluated more systematically.
    \item Continue using the \gls{adc} as part of regular deployment, with periodic review, further fine-tuning, and retraining cycles based on newly collected field samples.
    \item Strengthen the hyperparameter tuning process for both the \gls{yolo} and \gls{ocr} models by using broader and more systematic searches, so that future model updates are based on more stable and better-validated configurations.
    \item Automate the improvement cycle further by reducing manual steps between \gls{adc} export, annotation verification, training, and redeployment, so that future model updates require less manual coordination.
    \item Perform real robot trials in hospital-like and, when possible, actual hospital environments so that the perception pipeline can be validated under realistic operational constraints and new deployment data can be collected systematically.
    \item Replace the current Jetson microSD-based storage setup with an SSD-based configuration to improve robustness for logging, dataset export, and long deployment sessions.
    \item Reimplement runtime-critical \gls{ros} nodes in C++ where appropriate, since the current Python implementation is practical for development but likely leaves performance headroom on the embedded platform.
    \item Evaluate the pipeline on additional embedded platforms to determine how transferable the optimization choices are beyond the current \gls{jetson}-based setup.
    \item Integrate the perception pipeline with downstream robot action modules so that the 3D target estimates can be validated in full interaction tasks rather than perception-only experiments.
\end{enumerate}

\section{Final reflection}

This project provides a concrete foundation for autonomous interaction with
elevator infrastructure in hospital-like environments. More broadly, it shows
that robust embedded perception is not achieved by model selection alone. The
best results emerged from treating dataset design, model choice, runtime
architecture, and deployment data collection as one connected system. That
systems perspective is likely to remain essential as the CareBOTS platform evolves
toward more complete autonomous operation in real healthcare environments.
